%=====================================================================
% Front matter: title page, certificate, declaration, acknowledgement,
% abstract, contents, list of figures, list of tables.
%=====================================================================
\begin{titlepage}
\begin{tikzpicture}[remember picture,overlay]
  \draw[line width=2.5pt] ([xshift=0.6in,yshift=-0.6in]current page.north west)
    rectangle ([xshift=-0.6in,yshift=0.6in]current page.south east);
\end{tikzpicture}
\centering
\thispagestyle{empty}
\begin{spacing}{1.0}
\vspace*{0.02in}

\includegraphics[width=1.6in]{figures/pes_logo.png}\par
\vspace{0.2in}
{\large\itshape Dissertation on\par}
\vspace{0.5em}
{\Large\bfseries ``Tinker RL Lab:\\[0.35em]
Post-Training of Large Language Models''\par}
\vspace{0.4em}
{\normalsize A multi-framework training harness and studies of GRPO signal starvation\par}

\vfill

{\normalsize\itshape Submitted in partial fulfilment of the requirements for the award of the degree of\par}
\vspace{0.6em}
{\large\bfseries Master of Technology\par}
\vspace{0.2em}
{\normalsize\bfseries in\par}
\vspace{0.2em}
{\large\bfseries Data Science and Machine Learning\par}
\vspace{0.2em}
{\large\bfseries UE20CS972 -- Project Phase -- 2\par}

\vfill

{\normalsize\itshape Submitted by:\par}
\vspace{0.3em}
{\large\bfseries Arvind C R \qquad PES2PGE24DS140\par}

\vspace{0.5em}
{\normalsize\itshape Under the guidance of\par}
\vspace{0.2em}
{\large\bfseries Ramesh Prakash Guledgudd\par}

\vspace{0.6em}
{\normalsize\bfseries\color{pesmaroon} May 2026 -- October 2026\par}

\vfill

{\normalsize\bfseries DEPARTMENT OF COMPUTER SCIENCE AND ENGINEERING\par}
{\normalsize FACULTY OF ENGINEERING\par}
{\large\bfseries PES UNIVERSITY\par}
{\footnotesize (Established under Karnataka Act No.\ 16 of 2013)\\
Electronic City, Hosur Road, Bengaluru -- 560 100, Karnataka, India\par}
\end{spacing}
\end{titlepage}

%---------------------------------------------------------------------
\begin{titlepage}
\centering
\thispagestyle{empty}
\vspace*{0.1in}
{\Large\bfseries PES UNIVERSITY\par}
\vspace{0.2in}
{\large\bfseries FACULTY OF ENGINEERING\par}
{\normalsize DEPARTMENT OF COMPUTER SCIENCE AND ENGINEERING\par}
\vspace{0.35in}
{\Large\bfseries CERTIFICATE\par}
\vspace{0.5in}
\begin{spacing}{1.4}
\noindent This is to certify that the project report entitled
\textbf{``Tinker RL Lab: Post-Training of Large Language Models --- A
multi-framework training harness and studies of GRPO signal starvation''} is a
bonafide record of the project work carried out by
\textbf{Arvind C R} (SRN: \textbf{PES2PGE24DS140}), in partial fulfilment of
the requirements for the award of the degree of Master of Technology in
Data Science and Machine Learning of PES University, during the period
May 2026 -- October 2026. It is certified that all corrections/suggestions
indicated for internal assessment have been incorporated in the report.

\vspace{0.3in}
The candidate's individual contribution, the inherited Semester-3 group
infrastructure, and the boundary between them are stated in the Declaration
and in Chapter 1 \S1.6.
\end{spacing}
\vspace{0.4in}
\noindent\begin{tabular}{@{}p{2.4in}p{2.4in}@{}}
\dotfill & \dotfill \\
\textbf{Ramesh Prakash Guledgudd} & \textbf{Head of the Department} \\
Project Guide & Department of CSE \\
\end{tabular}
\vspace{0.3in}
\noindent\begin{tabular}{@{}p{2.4in}@{}}
\dotfill \\
\textbf{External Examiner} \\
\end{tabular}
\vfill
{\normalsize October 2026\par}
\end{titlepage}

%---------------------------------------------------------------------
\chapter*{Declaration}
\addcontentsline{toc}{chapter}{Declaration}
\thispagestyle{plain}

I declare that this project report is my own work. Where the work of others has
been used, it is cited and acknowledged. The matter embodied in this report has
not been submitted, in part or in full, to this or any other university or
institution for the award of any other degree or diploma.

The multi-framework benchmark infrastructure, the baseline framework
integrations, and the initial GRPO/PPO/DPO experimentation were produced by the
Semester-3 project group, of which I was a member, and are inherited here as
the shared foundation. The research direction, the experimental expansion, the
statistical analyses, the P1--P8 internal reports, the reporting and registry
artefacts, the adaptive group-size controller, the applied fraud study, and the
E1--E14 evaluation campaign reported in this document are my
individual contribution, carried out under the guidance of my project guide.

Every quantitative result in this report is bound to an evidence file in the
project repository. Where a result is incomplete, blocked, or negative, it is
reported as such rather than omitted. The boundary between inherited and
individual work is stated precisely in \texttt{platform\_hybrid/sem 4
work/PROVENANCE.md}.

\vspace{0.3in}
\noindent\begin{tabular}{@{}p{2.6in}@{}}
\includegraphics[height=0.55in]{figures/signature_arvind.png} \\
\dotfill \\
\textbf{Arvind C R} \\
SRN: PES2PGE24DS140 \\
\end{tabular}

%---------------------------------------------------------------------
% NEEDS AUTHOR CONFIRMATION: drafted from the tools recorded in the repository
% (git Co-Authored-By trailers, frontier_synthesis_*.tex, curriculum_opening/
% FINDINGS.md, zvf-program/jev_lab/, thesis_review_2026-09-27/). Confirm the
% list is complete and matches PES policy wording before submission.
\chapter*{Declaration on the Use of AI Tools}
\addcontentsline{toc}{chapter}{Declaration on the Use of AI Tools}
\thispagestyle{plain}

The following AI tools were used in this work, as recorded in the project
repository.

\begin{itemize}
\item \textbf{Coding and writing assistants.} Anthropic Claude models and
OpenAI Codex were used to write and revise code, analysis scripts and
documents, including drafts of this report. Commits they contributed carry a
\texttt{Co-Authored-By} trailer in the git history.
\item \textbf{Frontier-model synthesis.} ChatGPT Pro Extended and Gemini Deep
Think produced proposed derivations and reasoning recorded in
\texttt{platform\_hybrid/paper/sections/frontier\_synthesis\_*.tex}, which those
documents mark as not being experiments. Where this report uses such a result,
it states its own derivation (Chapter 6 \S6.3.3).
\item \textbf{Independent verification.} An LLM verification pass (Kimi,
6 July 2026) reviewed two single-seed results and flagged both as
\textsc{suspect}; both were then re-tested and reported as nulls (Chapter 6
\S6.4.3).
\item \textbf{Evaluation triage.} The Jev (TypeSafe) judgment harness was used
in the E1--E14 ledger only for classification, faithfulness and value
ordering, never for arithmetic (Chapter 4 \S4.5).
\item \textbf{Review.} The pre-defence self-review of 27 September 2026 and the
revision of this report that followed it were AI-assisted.
\end{itemize}

No AI assistant generated experimental data: every reported measurement comes
from an executed run or evaluation recorded in the repository, and simulated
values are labelled as simulations. The research direction,
the experimental design, the choice of claims and their final wording are the
author's responsibility.

%---------------------------------------------------------------------
\chapter*{Acknowledgement}
\addcontentsline{toc}{chapter}{Acknowledgement}
\thispagestyle{plain}

I thank my project guide, \textbf{Ramesh Prakash Guledgudd}, for direction and
sustained feedback throughout this work, and for the freedom to pursue the
conservative and negative results that form a significant part of it.

I thank the Semester-3 project group for the shared multi-framework benchmark
and the baseline integrations on which the Semester-4 studies build. I thank
the maintainers of the benchmark suites used in the E1--E14 campaign, whose
native evaluators made an independently checkable evaluation possible, and the
platforms --- Tinker, Modal, and Google Cloud --- on which the experiments ran.

Finally, I thank the reviewers whose questions at earlier reviews shaped the
reporting discipline this report depends on: the insistence on stating what a
number does \emph{not} establish has improved the work materially.

\vspace{0.3in}
\noindent\textbf{Arvind C R}

%---------------------------------------------------------------------
\chapter*{Abstract}
\addcontentsline{toc}{chapter}{Abstract}
\thispagestyle{plain}

This report studies group-relative reinforcement-learning post-training of
large language models under one methodological commitment: a comparative
claim is made only when the whole training stack is held fixed except the
factor under test. It makes four contributions.

\textbf{(1) A recomputable signal-starvation diagnostic and a loss-form
audit.} The Zero-Variance Fraction (ZVF) is logged per step and recomputed
from stored reward tensors, with a closed-form Bernoulli null. Reading the
runners' code rather than the internal reports shows that the ``GRPO'' runners
used here are REINFORCE with a group-mean baseline, without KL or clipping.

\textbf{(2) Stack-conditioned reporting.} A minimum reporting standard and a
machine-readable stack registry with a \texttt{stackdiff} tool that flags
label-flip risk between two runs from their records alone.

\textbf{(3) Calibrated negative results.} No reliable scaling law is recovered
from reward curves across twelve fitted anchors. A same-stack PPO/GRPO
contrast (0.992 against 0.990, $p = 0.374$, five seeds) and a group-size sweep
find no detectable difference, but both run on a saturated task and establish
neither equivalence nor an optimum. ZVF does not predict held-out accuracy
($\rho = 0.27$, 95\% CI $[-0.37, 0.88]$). An adaptive group-size controller
derived from ZVF saves no prompts in counterfactual simulation, because 92\% of its escalations fall on
groups that are already all-correct.

\textbf{(4) An evaluation-governance protocol.} One lightly trained actor
(Qwen3.6-35B-A3B with a LoRA adapter, 40 training steps) is evaluated on
fourteen held-out suites under native graders at pinned revisions, with
receipts, a no-pooling rule and a per-lane serving audit that traced a
zero score to a serving defect. Complete scopes include VerilogEval
($129/312 = 41.35\%$), LAB-Bench ($0.2288$) and AgentDojo benign utility
($88/97$); six replacement scopes completed on 27 September 2026
(Chapter~8, Table~8.C); seven original contracts remain externally blocked. A
paired trained-versus-base rerun on 5--151 identical items per lane finds no
significant difference in any lane, which at those sizes is inconclusive. No
claim of improvement over any baseline is made.

\vspace{0.25in}
\noindent\textbf{Keywords:} GRPO, RLVR, zero-variance fraction, signal
starvation, stack-conditioned results, benchmark evaluation, reproducibility,
held-out evaluation, attribution.

%---------------------------------------------------------------------
\tableofcontents
\listoffigures
\listoftables
